% !TEX root = ../../main.tex

\section{手势识别服务技术}

\subsection{FastAPI}

FastAPI 是一个基于 Python 的 Web 框架，适合构建高性能 API 服务，并能够自动生成交互式 API 文档\cite{fastapi-docs}。
在 HearBridge 系统中，Python 手势识别服务采用 FastAPI 作为接口服务框架，对外提供实时手势识别、样本采集、样本扫描、特征转换、模型训练、模型状态查询、模型重载和句子视频识别等能力。

将识别服务独立为 FastAPI 应用具有两个优点。
第一，Python 生态更适合集成 MediaPipe、TensorFlow、OpenCV 等视觉识别和模型训练工具，便于快速完成算法实验和服务封装。
第二，识别服务与 Spring Boot 业务服务解耦后，模型训练、模型加载和推理策略可以独立迭代，不会直接影响主业务服务的稳定性。

\subsection{MediaPipe}

MediaPipe 是 Google 提出的感知管线框架，可用于构建跨平台的视觉感知应用\cite{lugaresi-mediapipe-2019}。
其中，MediaPipe Hands 能够在设备端实现实时手部关键点跟踪，并输出手部骨架关键点信息\cite{zhang-mediapipe-hands-2020}。
在手势识别任务中，直接使用原始图像进行模型训练容易受到背景、光照、衣着、摄像头角度和设备差异影响；
而使用关键点特征可以将视频帧转换为结构化坐标序列，减少无关图像信息干扰。

在 HearBridge 识别服务中，MediaPipe 主要用于从移动端上传的视频帧或实时图像帧中提取手部关键点和人体姿态关键点。
系统基于这些关键点进一步构造手型、相对位置、角度和动态变化等特征，使模型能够关注手势动作本身。
对于固定词表下的手势训练任务，这种关键点特征方案比端到端原始视频模型更轻量，也更适合当前系统的开发周期和部署条件。

\subsection{TensorFlow/Keras}

TensorFlow 是 Google 开源的机器学习框架，支持模型构建、训练和部署，可用于图像识别、序列建模、自然语言处理等多种机器学习任务\cite{tensorflow-docs}。
Keras 是 TensorFlow 生态中常用的高层神经网络接口，能够以较简洁的方式搭建和训练深度学习模型。

在本系统中，TensorFlow/Keras 并不是直接对原始视频进行端到端识别，而是用于构建固定词表下的时序分类模型。
具体而言，Python 识别服务首先通过视频读取和关键点检测获得连续帧数据，再将关键点转换为固定长度的时序特征窗口，随后将该特征输入 TensorFlow/Keras 模型进行词级分类。
模型输出的类别标签、置信度和候选结果再经过后处理，用于实时识别结果展示、训练反馈或句子视频识别中的词序列生成。

\subsection{WebSocket}

WebSocket 适合实时双向通信，因此被用于移动端实时手势识别和样本采集场景。
移动端可以持续向 Python 识别服务发送图像帧，识别服务则持续返回当前动作识别结果或采集状态。
这种方式比频繁建立 HTTP 请求更适合实时相机帧传输和低延迟结果反馈。

在 HearBridge 系统中，WebSocket 主要解决实时识别过程中的连续帧传输问题。
移动端不断发送相机图像帧，识别服务在接收到足够的有效帧后进行窗口推理，并返回当前识别标签、置信度和模型状态。
通过 WebSocket，系统可以较自然地实现“边练习、边识别、边反馈”的训练体验。
